\section{引言}

海洋大数据的计算引擎正从岸基超算中心加速向海洋移动观测平台的边缘端延伸，使得“继续计算、补充观测还是立即行动”成为自主水下航行器必须面对的在线决策问题。传统的海洋数据处理依赖岸基中心汇集卫星遥感与数值模式输出，再将规划指令下发至水下平台\cite{chassignet2007hycom,chin2017mur}。然而，水下声学通信带宽极窄且上浮卫星通信代价高昂，航行器在执行长航程任务时必须依赖搭载的低功耗嵌入式计算模块独立完成环境推演与航路重规划\cite{remus100,jetsonorinnx}。由于机载电池总能量与算力极为有限，航行器既不能无限制地消耗算力进行高分辨率路径搜索，也不能频繁中断航行去执行耗时的原位测流或上浮通信\cite{horvitz1987reasoning,russell1991right}。因此，如何在计算、观测与行动之间合理分配受限的机载资源，构成了海洋边缘智能的核心挑战\cite{sung2021quit,magnino2026metareasoning}。

在台湾以东至东海的黑潮等强流海区，当自主水下航行器仅携带数天前下发的粗分辨率海流图执行绕障航行时，选错绕行方向将导致数小时的航时损失。长航程航行器的典型对水巡航速度约为1.0米/秒，而黑潮主轴及其邻近中尺度涡旋的流速可达到甚至超过这一量级。当航行器需要在相距百余千米的起终点之间绕越岛礁或禁航区时，从左侧还是右侧绕行是一个不可逆的二元拓扑决策：一旦驶入逆流一侧，不仅航时剧增，还可能因顶流过强而无法通行。由于水下无法实时接收高分辨率分析场，本文考虑航行器只携带下水前数天制作的1°粗分辨率海流图的情形：这张图既丢失了数十千米尺度的锋面与涡旋切变，又包含了数天内流轴摆动带来的时效误差。在这样粗糙且陈旧的地图下判断左右两侧孰快孰慢，极易因地图失真而误入慢侧。

面对粗图上的航路决策，最直接的常规做法是在当前最佳估计的海流图上投入更多机载算力，通过加密空间网格与航向离散度来求取更精确的最优航线\cite{lolla2014levelset}，但当环境信息本身缺失时，这种确定性等价的做法难以奏效。按确定性规划的原理，若左右两侧在粗网格上的预计航时十分接近，规划器会自然地将算力用于把网格进一步细分，以期分辨出哪一侧更快。然而，这一做法混淆了求解器的粗糙与海洋环境本身的不确定性：加密网格只能减少在同一张旧粗图上算路时的几何与离散误差，却不能为航行器提供任何关于真实涡旋位置的新信息\cite{jiang2015horizon,cuadron2025overthinking}。正如本文前期在东海赤潮巡测中的研究所示，当背景场存在明显时效误差时，成百上千倍地增加机载推演次数只会使规划器在自身想象模型上的打分不断虚高，而航线在真实海洋中的收益早已饱和\cite{smith2006curse}。

第二种常见路径是利用集合预报或随机流体模型开展风险感知航路规划，但现有方法不仅无法回答何时应当停止计算或去观测，还容易因有限集合抽样而漏掉真正翻转决策的流场\cite{wang2016path,subramani2018stochastic,kularatne2018forecast}。集合规划通过在数十个扰动流场中评估备选航路的平均航时或航时方差，试图选出期望最快或风险较小的路径。然而，这类方法将不确定性视为静态给定的输入，只输出一条折中航线，并未建立在算力消耗与观测代价之间进行切换的元推理机制\cite{lin2015metareasoning,elboher2023formal}。进一步而言，海洋流场的误差自由度极高，随机抽取的几十个集合成员往往扎堆于概率密度较高的平庸状态，难以覆盖合理误差球边界上那些虽然偏离中心、却恰好处在物理相容范围内并使所选航路急剧恶化的不利流场，从而低估了选错一侧的真实尾部风险\cite{benhaim2006infogap,marx2020predictive}。

第三种路径是基于信息价值的自适应采样与实验设计，即主动前往不确定性最高的海区进行测量，但对整个海流场有价值的信息并不等同于对当前左右绕行决策有价值的信息\cite{fossum2019,hollinger2014sampling}。传统的海洋自适应采样通常以降低全场估计方差或最大化互信息为目标，引导航行器前往误差最大的区域巡测\cite{srinivas2012gpucb}。然而，在绕障航行任务中，海流图在禁航区之外某处的巨大误差若对左右两侧航时的影响相同，或者该处根本不在备选航线上，测量它就纯粹是浪费航时；反之，顺路在起点附近测量虽然代价低廉，若决定两侧快慢的关键流场切变远在数十千米之外，近处测量也无法改变决策。序贯实验设计、模型判别与最优臂识别理论给出了正确的原则：真正决定是否需要新信息的，是在现有证据下能否找到一个与当前观测同样相容、却要求采取相反行动的竞争模型\cite{chernoff1959sequential,box1967discrimination,garivier2016optimal,attia2018goal}。

为此，本文提出“合理反例检验”这一核心思想：机载计算与现场观测消除的是两类性质截然不同的不确定性；当左右两侧航路难分伯仲仅仅是因为求解网格粗糙（即数值歧义）时，多算就能做对决定；而当存在与现有粗图同样相容、却使另一侧明显更快的合理海流场（即环境歧义）时，无论投入多少算力都无法分辨真伪，唯有引入新观测才能消解，且在观测代价过高时应当直接选择最坏合理损失更小的一侧。我们将这一判别转化为一个可直接求解的操作问题：针对粗图推荐的一侧航路，能否在由多年海洋大数据标定的合理海流集合内，主动搜出一个使另一侧反而更快、且让当前选择付出沉重航时代价的合理反例？图\ref{fig:overview}展示了从测试年份中按固定规则（照精算地图走选错且多花3小时以上、存在合理反例、门控后取保守选对的第一个任务）选出的2022年1月18日真实任务示例：在3天前的1°粗图上，左侧与右侧航时分别为37.71小时与38.02小时，地图因微弱优势推荐左侧；继续加密网格计算只会将这0.3小时的账面差值算得更精细，而合理反例检验则在历史误差标定的合理集合内搜出一个使左侧耗时55.69小时、右侧仅需44.07小时的物理相容流场，准确预警了当天真实海洋中左侧53.44小时、右侧41.72小时（右侧快11.7小时）的巨大逆流风险，并指导航行器根据最坏后悔执行门控决策。

\begin{figure}[htbp]
\centering
\includegraphics[width=\linewidth]{fig_route_overview}
\caption{粗分辨率海流图下的航路决策与合理反例检验（2022年1月18日任务示例）。(a) 航行器携带的3天前1°粗图，左侧航时37.7小时、右侧38.0小时，地图选左侧；(b) 在大数据标定的合理集合内搜出的合理反例海流，左侧55.7小时、右侧44.1小时；(c) 当天1/12°真实海流，左侧53.4小时、右侧41.7小时，右侧实际快11.7小时；(d) 本文决策流程：先校验粗图两侧时间差是否小于数值误差容限以决定是否按需精算，再通过反例检验计算最坏后悔，最坏后悔小则照地图走，大则在观测便宜时上浮取新图、观测昂贵时改选最坏后悔更小的保守一侧。再多的机载计算也只能把粗图上的0.3小时差值算准，唯有合理反例检验能显式揪出(b)(c)这类同样合理却翻转决策的海洋状态。}
\label{fig:overview}
\end{figure}

本文的主要贡献如下：
\begin{itemize}
    \item \textbf{概念与判别准则创新：}据我们所知，首次在自主水下航行器航路决策中严格分离了可由计算消除的“数值歧义”与只能靠新信息或稳健决策应对的“环境歧义”，提出了以“是否存在高损失合理反例”为核心判据的合理反例检验框架，明确了边缘计算与海洋观测的功能边界。
    \item \textbf{大数据标定与高效反例搜索方法：}提出了一种融合海洋大数据先验与包络定理梯度的反例检验算法；利用多年历史再分析海流标定旧粗图的无散流函数误差模态及覆盖90\%以上真实时间差的合理集合半径，将非凸流场搜索转化为球约束下的投影梯度上升，以相当于32成员集合的算力开销直接定位最坏后悔流场。
    \item \textbf{黑潮强流区系统实验与实证发现：}在六年黑潮再分析海流的1333个强流任务上证实，64倍算力仅修复28\%的粗算错误，而99.7\%的粗图决策存在合理反例；证明了决定行动的不是反例是否存在而是最坏后悔的大小（前20\%高风险任务集中67\%损失），所提门控保守策略将平均航时损失显著降低28\%，并量化了原位测流与上浮取图的收益边界。
\end{itemize}

本文其余部分组织如下：第2节回顾元推理、面向决策的信息获取及不确定海流路径规划的相关工作；第3节简述东海赤潮巡测前期研究，说明单纯增加计算为何收效甚微；第4节建立粗分辨率海流图下的绕障航路决策模型并剖析问题本质；第5节详述数值校验、大数据合理集合标定、投影梯度反例搜索与门控响应方法；第6节给出黑潮海区的实验结果、归因分析与适用范围；第7节拆解本文的原创创新点与研究演进脉络；第8节总结全文并展望未来工作。
